\section{Optimal path systems}\label{sec:paths}

We now assume $k\geq5$ and work in the graph $G$ from the
minimal-counterexample reduction. Thus $G$ has no cycle of order $k+1$,
$\delta(G)\geq k$, $\alpha(G)\leq k$, and the independent-set expansion
of Lemma~\ref{red:expansion} holds. Fix a maximum clique $Q$, and write
\[
 t=|Q|,\qquad h=k+1-t.
\]
By Theorem~\ref{clq:bound},
$\max(3,\lfloor k/2\rfloor)\leq t\leq k$, so $h\geq1$.
Our first objective is to encode paths through vertices outside $Q$ in a
way that allows us to close them into a cycle of exactly the forbidden
order. We then choose an optimal encoding and relate short paths outside
its vertex set to separated sets of vertices.

\subsection{Paths joining clique vertices}

A \emph{path system on $Q$} is a collection $\mathcal R$ of simple paths
with the following properties. Each path has two distinct endpoints in
$Q$ and a nonempty interior contained in $V(G)\setminus Q$. The interiors
of different paths are disjoint. Represent each path by the edge joining
its endpoints in an auxiliary graph with vertex set $Q$. These represented
edges must be distinct and must form a linear forest: every component is
a path, including isolated vertices. We call these components
\emph{chains}. In particular, every unused clique vertex is a singleton
chain; it does not represent a path with empty interior.

The \emph{amount} of an assigned path is its number of internal vertices.
For a system $\mathcal R$, let $q(\mathcal R)$ be the sum of its amounts,
let $e(\mathcal R)$ be its number of assigned paths, and let
$v(\mathcal R)$ be the number of clique vertices incident with its
represented edges. Thus singleton chains contribute neither to the amount
nor to the incident-vertex count. The empty collection is allowed.

\begin{lemma}[Closing a path system]\label{path:interval}
Let $G$ be a finite simple graph with no cycle of order $k+1$, and let
$Q$ be a clique of size $t$, where $3\leq t\leq k$.
Put $h=k+1-t$. No path system on $Q$, of total amount $q$ and with $v$
incident clique vertices, satisfies
\begin{equation}\label{path:forbidden-interval}
 h\leq q\leq k+1-v.
\end{equation}
\end{lemma}

\begin{proof}
Suppose such a system exists. The integer
\[
 u=k+1-q-v
\]
satisfies $0\leq u\leq t-v$: the first inequality follows from the upper
bound on $q$, and the second from $q\geq h$. Choose exactly $u$ of the
clique vertices not incident with represented edges.

Expand each nontrivial chain by replacing its represented edges with
their assigned paths. Each expansion is a simple path, and different
expansions are vertex-disjoint. Indeed, their clique vertices belong to
different forest components, and all assigned interiors are disjoint and
lie outside $Q$. Adjoin the $u$ chosen unused clique vertices as singleton
paths. These paths together contain exactly
\[
 q+v+u=k+1
\]
vertices, and all their endpoints belong to $Q$.

There is at least one nontrivial expanded chain, since $q\geq h\geq1$.
If there is only one path in the resulting list, it has at least three
vertices: it contains an assigned path with nonempty interior. The clique
edge between its distinct endpoints closes it into a simple cycle. If
there are two or more paths, order and orient them arbitrarily, join each
path's last endpoint to the next path's first endpoint by a clique edge,
and close the last to the first. This again gives a simple cycle. In the
case of two paths with one singleton, the two joining edges use the two
distinct endpoints of the nontrivial path; thus they are distinct. The
case of two singleton paths cannot occur. All remaining cases follow
directly from vertex-disjointness. The resulting cycle has order $k+1$,
a contradiction.
\end{proof}

Choose a path system $\mathcal P$ of maximum total amount $L$ subject to
$L<h$. Among systems with this amount, choose one with the minimum number
$e$ of assigned paths. This choice exists: the empty system has amount
zero, and $G$ is finite. Write
\begin{equation}\label{path:system-vertices}
 S=Q\cup\bigcup_{P\in\mathcal P}\bigl(V(P)\setminus Q\bigr),
 \qquad F=G-S.
\end{equation}
Since the amounts are positive integers and the interiors are disjoint,
\begin{equation}\label{path:budget}
 0\leq e\leq L\leq h-1=k-t,
 \qquad |S|=t+L\leq k.
\end{equation}
In particular, $L=0$ forces $e=0$; this includes the case $t=k$.

The choice of $\mathcal P$ and Lemma~\ref{path:interval} give a single
criterion that we will apply to modified systems.

\begin{lemma}[Optimality criterion]\label{path:optimality}
Let $\mathcal R$ be any path system on $Q$, with total amount $A$, with
$e'$ assigned paths, and with $v'$ incident clique vertices. Then
\begin{equation}\label{path:optimality-interval}
 L+\mathbf 1_{\{e'\geq e\}}\leq A\leq k+1-v'
\end{equation}
is impossible, where the indicator is one when $e'\geq e$ and zero
otherwise.
\end{lemma}

\begin{proof}
If $A\geq h$, Lemma~\ref{path:interval} applies. If $A<h$ and $e'\geq e$,
the lower bound gives $A\geq L+1$, contrary to the maximality of $L$.
If $A<h$ and $e'<e$, that lower bound gives $A\geq L$. Strict inequality
again contradicts maximality, and equality contradicts the choice of $e$.
\end{proof}

\begin{lemma}\label{path:exterior}
Every vertex of $S$ has at least $k+1-|S|\geq1$ neighbors in $F$.
\end{lemma}

\begin{proof}
A vertex of $S$ has at most $|S|-1$ neighbors in $S$ and has degree at
least $k$. Now use~\eqref{path:budget}.
\end{proof}

\subsection{Outside paths and separated balls}

The next definitions apply to any set $X\subseteq V(G)$; later $X$ will
be $S$ or a set obtained by adding vertices to $S$. For distinct
$x,y\in X$, an \emph{outside path relative to $X$ with parameter $d$}
is a simple $x$--$y$ path with exactly $d$ internal vertices, all in
$V(G)\setminus X$. Its length is $d+1$. We allow $d=0$, which means
the edge $xy$.

For $x\in X$, define subsets of $V(G)\setminus X$ by
\begin{equation}\label{path:balls}
 B_0(x;X)=N_G(x)\setminus X,
 \qquad
 B_{r+1}(x;X)=N_{G-X}[B_r(x;X)]\quad(r\geq0).
\end{equation}
Thus $B_r(x;X)$ consists of the vertices at distance at most $r$ in
$G-X$ from the outside-neighbor set $B_0(x;X)$. Every vertex of this
ball is joined to $x$ by a path of at most $r+1$ edges whose remaining
vertices lie outside $X$. The balls are nested and may be empty.

\begin{lemma}[Separation by forbidden outside paths]\label{path:separation}
Let $X\subseteq V(G)$ and let $x,y\in X$ be distinct.
\begin{enumerate}
\item For integers $r,s\geq0$, if no outside $x$--$y$ path relative to
$X$ has parameter $d$ with $1\leq d\leq r+s+2$, then
$B_r(x;X)$ and $B_s(y;X)$ are disjoint and anticomplete.
\item For an integer $r\geq1$, if no such path has parameter
$1\leq d\leq r$, then $y$ has no neighbor in $B_{r-1}(x;X)$.
\end{enumerate}
\end{lemma}

\begin{proof}
If the two balls in the first assertion meet, concatenate their paths
from $x$ and $y$ to a common vertex. This gives an $x$--$y$ walk of at
most $r+s+2$ edges, all of whose internal vertices lie outside $X$.
Deleting repeated portions produces a simple outside path with at most
$r+s+1$ internal vertices. If instead an edge joins the two balls, use
that edge between the two paths. The resulting walk has at most
$r+s+3$ edges and yields an outside path with at most $r+s+2$ internal
vertices. In both cases the parameter is at least one, since the walk
has no direct step from $x$ to $y$. Each case contradicts the hypothesis.

For the second assertion, a vertex of $B_{r-1}(x;X)$ has a path from
$x$ of at most $r$ edges with all subsequent vertices outside $X$.
If that vertex were adjacent to $y$, appending the edge to $y$ would
give an outside path with between one and $r$ internal vertices.
\end{proof}

When choosing separated sets, we also allow the singleton $\{x\}$ as
an option with \emph{radius label $-1$}. This is a separate convention;
it is not an additional step of the recursion~\eqref{path:balls}.
An option with radius $r\geq0$ means the ball $B_r(x;X)$.

\begin{corollary}[Singleton and ball compatibility]\label{path:singleton}
Consider two options based at distinct vertices $x,y\in X$, with
integer radius labels $r,s\geq-1$. If $r=s=-1$, the two options are disjoint
and anticomplete whenever $xy\notin E(G)$. Otherwise, they are disjoint
and anticomplete whenever all outside $x$--$y$ paths relative to $X$
with parameters
\[
 1\leq d\leq r+s+2
\]
are forbidden.
\end{corollary}

\begin{proof}
The case of two singletons is immediate, and the case $r,s\geq0$ is
Lemma~\ref{path:separation}. It remains to consider $r=-1$ and $s\geq0$,
after interchanging the options if necessary. The singleton $\{x\}$
and the ball $B_s(y;X)$ are disjoint because they lie in $X$ and
$V(G)\setminus X$, respectively. An edge between them would give an
outside $x$--$y$ path with between one and $s+1=r+s+2$ internal vertices,
contrary to the hypothesis.
\end{proof}

Pairwise separated options allow independent sets chosen in them to be
combined. In particular, if their individual independence numbers are
at least $u_1,\ldots,u_j$, then $G$ has an independent set of size at
least $u_1+\cdots+u_j$. We will use this observation to turn the path
restrictions imposed by optimality into a contradiction to
$\alpha(G)\leq k$.
